\begin{table}[h]
	\caption{Comparison between PEG and Pigeon}
	\label{table:peg_pigeon_comparison}
	\centering
	\resizebox{\textwidth}{!}{\begin{tabular}{cp{70mm}p{70mm}}
		& Pigeon & Peg \\\hline
		Installation & Requires running a \lstinline[language=Go]!go install! command and setting up environment variables & Requires supporting bash scripting (through Window's Subsystem for Linux (WSL) for Windows users), as well as running a \lstinline[language=Go]!go install! command, cloning the repository, and running \lstinline[language=Go]!go generate \&\& go build! \\
		Documentation & Contains a Godoc page, wiki page, an examples subdirectory, and a detailed README.md providing what PEG features are supported, command line interface usage, installation, and examples & Contains a detailed README.md and a PEG file syntax file containing what PEG features are supported, installation steps, and example projects links \\
		Generated Code & Generated Go code can be used to configure settings on the parser, view statistics or error information, debug, enable or disable memoization, and parse readers, files, and strings. Additionally, the generated code is well documented, with several comments surrounding usage of public functions & Generated Go code can be used to run the PEG script on a provided string, debug and output the Abstract Syntax Tree (AST), output errors in parsing, and enable or disable memoization \\
		Examples & Provides three direct examples: detecting indentation in files, parsing JSON files, and a calculator example. Examples were kept short (in regards to PEG scripts) with the calculator and indentation examples showcasing usage via a main function & Provides two project links for examples on PEG usage: a natural date/time parser and a scientific names parser. Both projects have large PEG scripts, with most lines being simple to read. The natural date/time project was intuitive to use and understand program flow \\
	\end{tabular}}
\end{table}